\chapter{INTRODUÇÃO}\label{chap:introducao}

A crescente diversidade das aplicações distribuídas modernas --- que abrangem desde transferências massivas de dados na Ciência de Dados Intensiva (\textit{Data-Intensive Science} -- DIS) até fluxos interativos sensíveis a atrasos --- impõe requisitos rigorosos e heterogêneos às infraestruturas de transporte \cite{paraiso2025dis, puttaswamy2025pathaware}. Para acomodar múltiplas demandas sobre uma mesma base física sem degradação mútua de desempenho, o fatiamento de redes (\textit{Network Slicing}) consolida-se como um paradigma essencial ao particionar os recursos de comutação e transmissão em redes lógicas isoladas, assegurando o cumprimento dos Acordos de Nível de Serviço (\textit{Service Level Agreements} -- SLAs) de cada classe de tráfego \cite{afolabi2018network, wei2026panqosr}.

Entretanto, a concretização de garantias granulares de Qualidade de Serviço (\textit{Quality of Service} -- QoS) exige que a arquitetura de rede ofereça visibilidade sobre o estado dos enlaces e controle explícito sobre as trajetórias percorridas pelos pacotes \cite{trammell2022rfc, borges2024pot}. Nesse contexto, emergem as Redes Cientes de Caminho (\textit{Path-Aware Networking} -- PAN), paradigma que supera o modelo tradicional de melhor esforço (\textit{best-effort}) ao expor atributos topológicos e de desempenho aos pontos finais ou aos dispositivos de borda, permitindo-lhes selecionar rotas otimizadas de acordo com os requisitos específicos de cada fluxo \cite{trammell2022rfc, wei2026panqosr}. Quando combinadas à programabilidade do plano de dados e a mecanismos de escalonamento diferenciado nas interfaces de saída, as redes cientes de caminho viabilizam a coexistência eficiente entre fluxos de longa duração e alta vazão (\textit{elephant flows}) e fluxos curtos críticos em latência (\textit{mice flows}) \cite{santos2024mtspolka, puttaswamy2025pathaware}.

\section{PROBLEMA}

Na arquitetura tradicional da Internet, o encaminhamento de pacotes orienta-se exclusivamente pelo endereço de destino por meio de decisões locais salto a salto (\textit{hop-by-hop}) baseadas em tabelas de roteamento distribuídas \cite{postel1981, borges2024pot}. Esse modelo comporta-se como uma infraestrutura opaca: os sistemas finais e os comutadores de borda desconhecem as propriedades das rotas subjacentes, não conseguem influenciar a seleção de caminhos alternativos e tampouco dispõem de garantias nativas de que os pacotes transitaram pela sequência de nós esperada \cite{trammell2022rfc, Unveiling_the_Mystery, borges2024pot}. Como consequência, protocolos convencionais tendem a concentrar o tráfego nos caminhos de menor contagem de saltos, gerando congestionamento em enlaces específicos enquanto rotas alternativas permanecem subutilizadas \cite{paraiso2025dis, mariano2026ml}.

Embora o paradigma de Redes Definidas por Software (\textit{Software-Defined Networking} -- SDN) e a introdução de planos de dados programáveis via linguagem P4 tenham conferido maior flexibilidade ao controle da rede \cite{KAUR2021109}, a implementação de fatiamento com garantias estritas de QoS ainda enfrenta barreiras significativas de escalabilidade e agilidade operacional:

\begin{itemize}
    \item \textbf{Sobrecarga de estado e lentidão de convergência:} Soluções tradicionais de engenharia de tráfego e reserva de recursos (como MPLS-TE, IntServ ou tabelas de grupos ponderados) exigem a manutenção de entradas por fluxo ou por fatia em todos os dispositivos intermediários \cite{wei2026panqosr, santos2024mtspolka}. A cada alteração de rota ou política de prioridade, o controlador precisa atualizar coordenadamente múltiplos comutadores no núcleo, elevando o tempo de convergência e a sobrecarga do plano de controle \cite{liberato2018rdna, froes2020proglab}.
    \item \textbf{Limitações do roteamento na origem baseado em listas:} Protocolos de roteamento na origem como o \textit{Segment Routing} (SR) mitigam o estado no núcleo ao embutir uma lista de identificadores de segmento no pacote \cite{rfc8402, ventre2020segment}. Contudo, essa abordagem exige operações sucessivas de leitura e atualização ou remoção de cabeçalhos a cada salto (\textit{label stack rewriting}), além de incorrer em sobrecarga variável proporcional ao comprimento do caminho \cite{dominicini2020, borges2024pot}.
    \item \textbf{Insuficiência do isolamento exclusivo por seleção de caminho:} Conforme evidenciado por estudos recentes em redes cientes de caminho \cite{wei2026panqosr, puttaswamy2025pathaware}, direcionar um fluxo para uma rota específica (\textit{path steering}) não basta para assegurar o cumprimento de SLAs quando múltiplas fatias compartilham os mesmos enlaces físicos. Sem mecanismos programáveis de enfileiramento e escalonamento ponderado no plano de dados, fluxos volumosos podem saturar os buffers intermediários e degradar severamente a latência de fluxos sensíveis ao atraso \cite{santos2024mtspolka, Elbediwy2025}.
\end{itemize}

A arquitetura PolKA (\textit{Polynomial Key-based Architecture}) \cite{dominicini2020} superou parte dessas limitações ao empregar o Sistema Numérico de Resíduos (\textit{Residue Number System} -- RNS) e aritmética polinomial sobre o corpo finito $GF(2)$ para codificar toda a rota em um único rótulo fixo e imutável na borda da rede, permitindo o encaminhamento totalmente sem estados (\textit{stateless}) no núcleo mediante o reuso de unidades de \textit{CRC} em hardware \cite{dominicini2021}. Evoluções subsequentes do ecossistema \textit{PolKA} demonstraram sua versatilidade em redes cientes de caminho para telemetria em banda (\textit{MPolKA-INT}) \cite{pereira2024mpolkaint}, comprovação de trânsito (\textit{PoT-PolKA}) \cite{borges2024pot}, divisão de tráfego multicaminho (\textit{MTS-PolKA}) \cite{santos2024mtspolka} e orquestração orientada a dados e aprendizado de máquina \cite{paraiso2025dis, mariano2026ml}. Entretanto, o protocolo PolKA original restringe-se à seleção da interface física de saída em cada salto, carecendo de primitivas nativas para codificar políticas de filas de prioridade e prover isolamento de recursos entre fatias de rede concorrentes.

\section{PROPOSTA}

Para preencher essa lacuna, esta dissertação propõe e avalia a arquitetura \nome\ (\textit{Slice-aware Polynomial Key-based Architecture}), uma extensão do roteamento na origem \textit{PolKA}, concebida para habilitar o fatiamento de redes e o provisionamento granular de QoS em redes cientes de caminho com núcleo totalmente sem estados.

A proposta unifica, em um único identificador polinomial de comprimento fixo inserido exclusivamente nos comutadores de borda, tanto a trajetória explícita pela topologia quanto a política de enfileiramento a ser aplicada em cada nó intermediário. No plano de dados dos comutadores de núcleo, cada dispositivo executa uma operação de módulo polinomial em $GF(2)$ entre o rótulo do pacote e seu identificador local, extraindo simultaneamente —-- por meio de fatiamento de bits (\textit{bit-slicing}) —-- a porta física de egresso e o identificador da fila de serviço. Ao direcionar os fluxos para filas previamente configuradas sob políticas de escalonamento ponderado nos nós intermediários, o \nome\ viabiliza o isolamento de largura de banda entre fatias concorrentes sem manter estado no núcleo e permite a migração instantânea de rotas e classes de serviço mediante a atualização exclusiva das regras na borda da rede.

\section{OBJETIVOS}

\subsection{Objetivo Geral}

Propor, implementar e avaliar o protocolo \nome, uma arquitetura de roteamento na origem sem estados no núcleo baseada em aritmética polinomial em $GF(2)$, voltada ao fatiamento de redes e ao provisionamento de garantias de Qualidade de Serviço (QoS) em redes cientes de caminho com planos de dados programáveis.

\subsection{Objetivos Específicos}

Para alcançar o objetivo geral, estabelecem-se os seguintes objetivos específicos:

\begin{enumerate}
    \item Investigar os fundamentos e os requisitos arquiteturais para o fatiamento de redes e o provimento de QoS granular em Redes Cientes de Caminho (\textit{Path-Aware Networks});
    \item Analisar as limitações do roteamento na origem baseado em \textit{PolKA} e dos mecanismos tradicionais de escalonamento em planos de dados programáveis quanto ao isolamento de tráfego entre fatias;
    \item Projetar a extensão polinomial e o mecanismo de decodificação simultânea de porta de saída e fila de serviço do protocolo \nome, mantendo o encaminhamento sem estados no núcleo da rede;
    \item Desenvolver e validar protótipos funcionais dos planos de dados e de controle do roteamento baseado em RNS em dois ambientes experimentais complementares: (i) um ambiente emulado em P4/BMv2 integrado a um gerenciador de tráfego customizado com escalonamento ponderado (Weighted Deficit Round-Robin – WDRR); e (ii) uma implementação de alto desempenho em contêineres Ubuntu 24.04 orquestrados via Containerlab, operando sobre programas eBPF e utilizando o subsistema nativo de controle de tráfego (Traffic Control – tc) do kernel Linux para o enfileiramento e escalonamento de pacotes.
    \item Avaliar experimentalmente o desempenho da arquitetura proposta em cenários com fluxos concorrentes, analisando métricas de isolamento de largura de banda, latência, vazão e agilidade na migração dinâmica de rotas e políticas de filas.
\end{enumerate}

\section{ORGANIZAÇÃO DO TRABALHO}

O restante desta dissertação está organizado da seguinte forma. O Capítulo~\ref{chap:referencialTeorico} apresenta o referencial teórico, abordando Redes Definidas por Software, programabilidade do plano de dados com P4 e eBPF, protocolos de roteamento na origem com ênfase na arquitetura PolKA, conceitos de fatiamento de redes e os fundamentos de Redes Cientes de Caminho. O Capítulo~\ref{chap:proposta} detalha os requisitos, a formulação matemática em $GF(2)$, o formato do cabeçalho, a decodificação no plano de dados e o gerenciador de tráfego programável da arquitetura \nome. O Capítulo~\ref{chap:prototipos} descreve a implementação dos protótipos desenvolvidos, a metodologia experimental e a análise dos resultados obtidos. Por fim, o Capítulo~\ref{chap:conclusao} apresenta as conclusões do trabalho e aponta direções para pesquisas futuras.